% Compile from this directory with: pdflatex project_14.tex (twice).
\documentclass[11pt,a4paper]{article}
\usepackage[T1]{fontenc}
\usepackage[utf8]{inputenc}
\usepackage[margin=22mm]{geometry}
\usepackage{lmodern}
\usepackage{amsmath,amssymb}
\usepackage{enumitem}
\usepackage{xcolor}
\usepackage{microtype}
\usepackage{hyperref}
\definecolor{epflred}{HTML}{E3000B}
\hypersetup{colorlinks=true,urlcolor=epflred,linkcolor=epflred,
  pdftitle={ENG-654 Project 14: Contour Placement for PUMA 560}}
\setlength{\parindent}{0pt}
\setlength{\parskip}{6pt}
\setlist[itemize]{leftmargin=*,itemsep=5pt,topsep=4pt}
\urlstyle{tt}

\begin{document}
{\small\textbf{EPFL \textbar\ ENG-654}\hfill Kinematics-Grounded Motion Planning for Robots}
\vspace{5mm}

{\color{epflred}\Large\bfseries Project 14}

{\LARGE\bfseries Contour Placement for PUMA 560\par}

\textbf{Format:} Group of two students; simulation-based project.\\
\textbf{Course foundations:} Lectures 4--7; Exercises 2--3.

\section*{Motivation}
A contour may be geometrically small and entirely reachable yet fail
from every admissible initial posture at a particular placement. Translating
or rotating the task changes the IK trajectories and their margins. This
project connects task placement to the global branch structure of a specific
noncuspidal industrial robot.

\section*{Task}
\textbf{Overall objectives.} Map where a fixed full-pose contour can be executed by PUMA 560, identify
which initial shoulder--elbow--wrist classes survive, and explain feasibility
boundaries through singularities and joint limits.

\begin{itemize}
  \item \textbf{Validate the PUMA full-pose model.} Reuse the course D--H model,
analytical IK, and signed shoulder--elbow--wrist classification. Verify
forward pose closure and document joint limits and base/tool transforms.
Keep all mathematical branches distinct before physical filtering.

  \item \textbf{Define a contour and placement family.} Specify one closed contour
with a smooth or explicitly segmented orientation schedule. Place the complete
pose sequence using $T'_k=T_\Delta T_k$, including orientation. Explore two
translation coordinates on a small grid and compare two fixed orientations
of the whole contour.

  \item \textbf{Continue every admissible initial class.} For each placement,
enumerate initial IKs and lift the full-pose path from every valid one.
Track class identity, position/orientation residuals, joint limits, and arm
and wrist regularity. Never repair a failed branch by jumping to another
configuration at the next sample.

  \item \textbf{Produce branch-specific placement maps.} For each initial class,
plot complete-path feasibility against the two placement coordinates.
Distinguish unreachable poses, branch termination, joint-limit violations,
and singularity-margin failures. Compare surviving trajectories by joint
travel and minimum margins.

  \item \textbf{Explain the boundaries geometrically.} Select two neighboring
placements with different outcomes. Plot the PUMA arm determinant in
$(q_2,q_3)$, overlay the wrist-center IK trajectory, and relate its workspace
image to arm IK-count regions. Diagnose wrist failures using the full
Jacobian rather than the reduced position plot alone.
\end{itemize}

\newpage
\section*{Specifics and scope}
Use one contour, a $5\times5$ translation grid, and two fixed placement
rotations; locally refine only one boundary. Begin with about 150 samples per
contour. Choose a justified characteristic length for the mixed pose Jacobian
and use the same thresholds throughout. This is kinematic placement analysis;
cycle-time optimization and collision-checker development are not required.

Check both translation and orientation residuals
for every claimed full-pose IK. Document angle periodicity and limits,
Jacobian conventions, and any scaling used for singular values. State all
fixed coordinates in reduced plots; a two-dimensional slice does not show
the entire 6R singular set. Refine decisive transitions, and distinguish
sampled evidence from a continuous-path guarantee. Collision freedom is
not implied by these kinematic checks.

\section*{Expected outcome}
Understand how task placement changes the feasible global IK choices
for PUMA. Deliver the contour and placement transforms, eight class-labelled
feasibility maps per placement orientation, and two detailed explanations
connecting placement boundaries to joint-space and workspace structure.

Submit reproducible code, model parameters and test data, the required plots,
a concise 5--6-page report, and a short simulation demonstration. Explain
both successful cases and failures; conclusions must follow the numerical
and geometric evidence rather than an expected outcome.

\section*{Resources}
The PUMA 560 URDF is
\path{lectures_main/assets/models/puma/puma560_robot.urdf}; its STL meshes are
in the same directory. The lecture website provides axis/frame inspection,
IK visualization, and path demonstrations. Confirm joint offsets and the
selected terminal frame when importing the model.
Reuse Lectures 4 and 6 for full-pose IK/classification and Lecture 7 for
continuation. Its existing PUMA path demo constrains wrist-center position
only; extend or complement it with the provided full-pose IK for this project.

Use the lecture website to verify D--H parameters, visualize IKs and
singularities, and simulate paths. All listed paths are relative to the
repository root. Start the website following \path{lectures_main/README.md}.

\section*{Workload and collaboration}
Budget approximately 60 hours per student (120 person-hours per pair)
for this 2-ECTS project, following EPFL's 30-hours-per-credit convention.\footnote{\href{https://www.epfl.ch/education/teaching/teaching-guide-2/getting-started/design-a-course_1/}{EPFL: Design a Course, workload guidance}.}
Deduct any lectures or exercises included in those credits. Allocate roughly
20 team-hours to model validation, 50 to the core work, 30 to experiments,
and 20 to reporting. Reuse the specified IK and simulations. Share validation
and interpretation even if modelling and implementation are divided between
students. Hardware execution is outside scope.

\end{document}
